Positive unit costs and Assumption~\(\mathrm{ass{:}nondegenerate\text{-}components}\) imply, for every \(d\in\mathcal D\),
\[
 0<c_{E,d}V_{E,d}<\infty,
 \qquad 0<c_{O,d}V_{O,d}<\infty.
\]
Theorem~\(\mathrm{thm{:}budget\text{-}profile}\) therefore gives
\[
 v_d:=\mathcal V_d^*(P)
 =(\sqrt{c_{E,d}V_{E,d}}+\sqrt{c_{O,d}V_{O,d}})^2\in(0,\infty).               \tag{5}
\]
For either \(j=0\) or \(j=A\), Definition~\(\mathrm{def{:}predictive\text{-}benchmark}\) gives
\[
 \rho_{\mathrm{pred}}^j(P)
 =\frac{\min_{d\in\mathcal D_{\mathrm{pred}}^j(P)}v_d}
 {\min_{d'\in\mathcal D}v_{d'}}.                                             \tag{6}
\]
Substitution of (5) proves (1).  The numerator minimizes over a nonempty subset of the denominator's index set, so the ratio is at least one; positivity and finiteness make it finite.  Dividing the best-tie inequality by the positive denominator in (6) shows that its least constant is exactly \(\rho_{\mathrm{pred}}^j(P)\).

The index set in (2) contains every diagonal pair and is finite.  Any constant valid for all implications must dominate every displayed ratio, while their maximum makes all implications valid.  This proves both exactness and lawwise finiteness of (2) for each risk rule.

Theorem~\(\mathrm{thm{:}two\text{-}rule\text{-}predictive\text{-}impossibility}\) proves (3), with paired center costs \(5\chi\) and \(5\), and Lemma~\(\mathrm{lem{:}sharp\text{-}uniform\text{-}binary\text{-}two\text{-}risk\text{-}elbow}\) proves the common optimizing sequence and (4).  These results preserve all four exact center ratios, so the conclusions hold for the unadjusted and treatment-aware rules under both normalizations.  This completes the proof.